\section{Introduction}
\label{sec:intro}

\subsection{The question}

Start from a flat real vector space $V=\R^{3,3}$ of split signature, with an orthonormal basis $e_1,\dots,e_6$ in which
$e_1,e_2,e_3$ have square $+1$ and $e_4,e_5,e_6$ have square $-1$, and with the map $e_i\mapsto e_{i+3}$ that pairs each
positive direction with a negative one. Read the positive triple as three times and the negative triple as three
lengths. The space has six directions, the same number as an electromagnetic field strength, and the pairing looks like
the correspondence between a magnetic and an electric component. This geometric picture, which we call the
\emph{electromagnetic sphere}, is the starting point of the paper.

The hypothesis we test is that as much of relativistic field theory as possible should be read off the internal
structure of the Clifford algebra $\Cl(3,3)\cong M_8(\R)$ and of its graded tower $\Cl(3n,3n)\cong M_{2^{3n}}(\R)$,
rather than added by hand. Split signature is the smallest setting in which three positive directions can be paired with
three negative ones by an anti-isometry, and the tower is the natural way to enlarge the arena without changing its
local structure. The classification of real Clifford algebras and their gradings is classical
\cite{LawsonMichelsohn,Lounesto,Porteous,MosnaMirallesVaz}, and several programmes have tried to extract Standard-Model
structure from division and Clifford algebras \cite{Furey,Gresnigt,ChamseddineConnes}; the representation-theoretic limits
of such embeddings are discussed in \cite{DistlerGaribaldi}.

The question we ask is not whether the Standard Model fits inside this arena. It does, in more than one way, and fitting
proves nothing. The question is which of the free choices of relativistic field theory the algebra removes, and which it
leaves free. That question has sharp answers, and the answers are the content of this paper.

\subsection{A deflationary principle}
\label{sec:deflation}

One theorem, proved in \cref{sec:dirac}, governs everything else. Take the four interior gamma matrices
$\gamma^\mu=(e_1,e_4,e_5,e_6)$ and the operator $J=e_2e_3$, which commutes with all four and squares to $-\id$. Then the
algebra generated over $\C$, with $J$ playing the role of $i$, by the four gamma matrices is the whole of $M_4(\C)$. The
real Dirac theory built on $\Cl(3,3)$ is therefore ordinary complex Dirac theory written in real variables, exactly and
not in a limit. A change of variables predicts nothing, so a statement about the arena carries information only if it
could have come out otherwise. We apply this test throughout. Facts that are true of Dirac theory in any notation, such as
$g=2$ at tree level, current conservation or the Ward identity, are \emph{inherited}; they are reproduced here only where
they are needed to state something else. The statements we report are of a different kind: which bilinears are admissible,
which generators may be gauged, how large the internal centraliser is at each rung, what the faithful gauge algebra of a
projected block is, which potentials can select a vacuum, and which constructions are impossible under stated
hypotheses.

Every claim in this paper belongs to one of four classes, and we say in the text which one. A result is \emph{fixed by the
algebra} when it is an identity or a theorem that could have failed. It is \emph{constructed} when the algebra admits it
and we exhibit it, without anything selecting it among alternatives. It is \emph{obstructed} when it is impossible under
hypotheses that we state, and those hypotheses are part of the result. It is \emph{open} when it is neither established
nor refuted. The finite-dimensional statements are exact: every blade of the realisation used here is a signed
permutation matrix, so products, commutators, traces and ranks are integer computations. How the statements were verified
is described in \cref{app:methods}.

\subsection{Summary of results}

\paragraph{The arena (\cref{sec:arena,sec:clock,sec:dirac,sec:discrete}).} A bivector $e_ae_b$ squares to $-\id$ when its
two directions have the same sign and to $+\id$ when they have opposite signs; the first kind generates compact rotations
that can be gauged, the second kind generates boosts that can only grade. The three \emph{pole operators}
$e_ie_{i+3}$ are commuting involutions grading the algebra by $\Z_2^3$, and they span a maximally split Cartan
subalgebra of $\so{3,3}$. The pairing has two square roots: a para-complex structure $\cK$ with $\cK^2=+\id$, which grades,
and a complex structure $G$ with $G^2=-\id$, which is self-adjoint for $\eta$ and, as we show later, binds the interior
dynamics and acts as a Hodge star. A clock is a unit vector of the positive triple; the symmetric clock
$\widehat W=(e_1+e_2+e_3)/\sqrt3$ is the unique permutation-invariant point and the entropy maximum of the clock simplex
under the Fisher--Rao metric, and exactly that clock makes the conversion operator invisible in the equations of motion.
The imaginary unit is selected rather than adjoined: inside the joint commutant of the interior Lorentz action and the
internal bivectors, the only square roots of $-\id$ that commute with the gamma matrices are $\pm J$. Charge conjugation is
the surplus plane $\langle e_2,e_3\rangle$ and coincides with the Majorana structure; time reversal is forced to square to
$-\id$; the kinetic forms alone fix the massless gauge algebra $\uu1\oplus\su2$, which the mass form cuts to the
electromagnetic $\uu1$; and the quantum state space must be formed by complexifying from outside.

\paragraph{The tower and its matter (\cref{sec:tower,sec:projectors,sec:pstar,sec:charges,sec:actions}).} The bivectors
commuting with the interior Lorentz action of $\Cl(3n,3n)$ form $\so{3n-1,3n-3}$, nested along the tower, of dimensions
$1,28,91,190,\dots$. At one constituent this centraliser is the single bivector $J$, so the photon is not one symmetry of the
interior structure but the only one. At three constituents we construct an orthogonal projector
\[
P_*=\tfrac18\,(\id+\Gamma\chi_W)(\id-J\chi_C)(\id-Jq_A)
\]
of rank $64$ which commutes with the Lorentz action and with $J$, is axial in the sense that $P_*\gamma^\mu P_*=0$, keeps the whole of
$\su4$ acting faithfully, and carries $(\mathbf4,\mathbf2,\mathbf1)_L\oplus(\mathbf4,\mathbf1,\mathbf2)_R$: one Pati--Salam
family with a right-handed neutrino. Its two central abelian factors cannot be gauged on this matter. Charges built from
three commuting involutions are conjugately paired whenever they are anomaly free, so hypercharge cannot be read off from
anomaly cancellation inside that class. When the Lorentz action is instead taken diagonally from all three constituents,
the internal algebra becomes the entirely compact $\so6\oplus\so3$; the copy algebra and the Lorentz algebra form a
reductive dual pair satisfying $C_{\rm cop}+C_\Delta=\tfrac{15}2$, and the module is vector-like. A single theorem gives
the candidate gauge algebra of every Lorentz action on the tower, and with it a rung-by-rung census of the vector bosons
the tower can carry.

\paragraph{Gauge fields and the electromagnetic reading (\cref{sec:fields,sec:emfield,sec:spin,sec:orth}).} A varying
projector induces a connection whose available curvature algebra is far too large, so a parent action would have to cut
the projector's freedom rather than build a gauge field. The centraliser of the complex structure $G$ in $\so{3,3}$ is a
Lorentz algebra, under which $\R^{3,3}$ \emph{is} $\Lambda^2\R^{1,3}$: boosts act by $\delta\bE=\bv\times\bB$,
$\delta\bB=-\bv\times\bE$, $G$ is the Hodge star, and the two invariants $|\bB|^2-|\bE|^2$ and $\bE\cdot\bB$ are the real
and imaginary parts of one holomorphic quadratic form. This reading excludes the interior four-plane on which the matter
is built. Covariance fixes the form of Maxwell's equations, because the vector module occurs in $\Lambda^2\otimes V$ with
multiplicity exactly two, and fixes nothing else. The photon's spin one appears in two places, related by the Klein
correspondence $V\cong\Lambda^2(S^+)^*$, and at the second rung both can be coupled to a Dirac block in a single gauge
theory. The orthogonality of $\bE$ and $\bB$ is the reality of the invariant $q=z\cdot z$; its failure breaks parity and time
reversal, is never spontaneous, and does not generate mass.

\paragraph{Interior dynamics (\cref{sec:shifted,sec:neutral,sec:closure}).} The Lagrangian of a single shifted operator
$\partial_\tau+\kappa g$ factorises exactly when the shift is symmetric, and then it describes an inverted oscillator; for
every constant shift some solution is unbounded, and with a positive form no solution with nonzero momentum recurs. In the
split form the same factorisation admits bound motion, and it does so exactly in neutral signature, with the complex
structure $G$ as the shift. The force law is a dual transport which is a Lie--Poisson equation only on an extension of the
frame; that extension is unique, its rotor angle is the duality angle of the momentum, and its free motion is a pendulum
that drifts the frame except on null data.

\paragraph{The vacuum (\cref{sec:subfields,sec:kernels,sec:selection,sec:matter,sec:goldstone}).} At one constituent we
classify the vacuum of every coordinate three-plane in four senses. At three constituents the space of Lorentz-invariant
mass kernels of the chiral block has dimension $272$ and decomposes as
$(\mathbf{10},\mathbf3,\mathbf1)\oplus(\mathbf{10},\mathbf1,\mathbf3)\oplus(\mathbf{15},\mathbf2,\mathbf2)\oplus
2(\mathbf6,\mathbf1,\mathbf1)\oplus(\mathbf1,\mathbf2,\mathbf2)$; the Standard-Model invariant kernels form one complex
line, the right-handed neutrino Majorana kernel. We give the complete renormalizable potential of the
$(\mathbf{10},\mathbf1,\mathbf3)$ component, including a holomorphic quartic that violates baryon number by two units,
and the necessary and sufficient conditions under which the Standard-Model orbit is its global minimum; we compute the
radiatively generated spectrum of all $271$ angular modes, the Goldstone content, the discrete symmetries of the vacuum,
and an electroweak stage with the Georgi--Jarlskog factor and a seesaw.

\paragraph{Curvature and the photon (\cref{sec:curved,sec:curvop,sec:transport,sec:photon}).} In a curved arena the
pairing becomes an almost para-Hermitian structure. An explicit six-dimensional metric with a normal $U(1)$ has a
closed-form Einstein tensor whose tangential block on the four-dimensional slice is the Einstein tensor of the slice, while
the normal field strength gravitates with negative energy, as the signature forces. In the field-strength reading the
curvature operator of a four-metric lives on the frame, the first Bianchi identity is the single condition
$\Tr(G\cR)=0$, and the duality rotation of the curvature ends at the double dual, which contracts to the Einstein tensor.
The projector current sources a massive normal vector, and a massless photon requires a vacuum that commutes with the
charge; such a vacuum exists at the second rung and not at the first.

\paragraph{Causal structure (\cref{sec:baseline,sec:selector,sec:typesel,sec:constraints}).} A field propagating freely
in six dimensions of signature $(3,3)$ is ill-posed and has no ground state, and its time-like cone has no past and future.
Choosing a Lorentzian four-plane inside $\R^{3,3}$ is the same as choosing a complex structure, and at one constituent the
light cone of that plane is the set of chiral spinors squared with the help of $J$. No bounded-below invariant potential
of degree below eight isolates that choice, and every invariant kinetic term for it has ghost directions. The surplus-time
constraints, by contrast, follow from a gauge symmetry of the worldline.

\subsection{Conventions}
\label{sec:conventions}

Spacetime signature is $(+,-,-,-)$, units are Heaviside--Lorentz with $\hbar=c=1$ unless stated. A \emph{constituent} is
a tensor factor $\Cl(3,3)$ of the tower; the word never refers to a particle, and three constituents do not describe three
particles. Anomaly coefficients are chiral traces per Weyl component with every field written left-handed. In the
field-strength reading of $\R^{3,3}$ the vector $X=\sum_i(B_ie_i+E_ie_{i+3})$ is identified with the two-form $F$ of
components $F_{0i}=E_i$, $F_{jk}=\epsilon_{ijk}B_i$; with this orientation $\bB$ is minus the textbook magnetic field,
which we write $\Btext=-\bB$, while $\bE$ is the textbook electric field. Formulas taken from the literature are written in
textbook variables, and where a sign matters the variables are named. Compact Lie algebras are realised by skew operators
and the corresponding charges by the symmetric operators $\hat t=-Jt$.

\subsection{Organisation}

\Cref{sec:arena,sec:clock,sec:dirac,sec:discrete} develop one constituent. \Cref{sec:tower,sec:projectors,sec:pstar,sec:charges,sec:actions} climb the tower and build the matter. \Cref{sec:fields,sec:emfield,sec:spin,sec:orth} treat gauge
fields and the field-strength reading, \cref{sec:shifted,sec:neutral,sec:closure} the interior dynamics, and
\cref{sec:subfields,sec:kernels,sec:selection,sec:matter,sec:goldstone} the vacuum. \Cref{sec:curved,sec:curvop,sec:transport,sec:photon} pass to a curved arena and to the photon, \cref{sec:baseline,sec:selector,sec:typesel,sec:constraints} to causal structure, and \cref{sec:assessment} collects what is fixed, what is constructed, what is
obstructed and what is open. \Cref{app:realisation} gives the explicit matrices, so that every finite-dimensional claim
can be checked by hand, and \cref{app:methods} describes the verification.
